\documentclass[11pt]{article}
\usepackage[T1]{fontenc}
\usepackage{lmodern,amsmath,amssymb,amsthm}
\usepackage[margin=1in]{geometry}
\usepackage[hidelinks]{hyperref}
\newtheorem{theorem}{Theorem}
\newtheorem{corollary}[theorem]{Corollary}
\newcommand{\Zp}{\mathbb Z_p}
\newcommand{\Qp}{\mathbb Q_p}
\newcommand{\OO}{\mathcal O}
\newcommand{\PP}{\mathbb P^1}
\newcommand{\adj}{\operatorname{adj}}
\newcommand{\cont}{\operatorname{cont}}
\title{Matrix Coset Averages in the Boundary Model of the Steinberg Representation}
\author{\small\href{mailto:admin@cosmosgeek.org}{admin@cosmosgeek.org}}
\date{Working draft for specialist review}
\begin{document}
\maketitle
\begin{abstract}
We compute the additive average of the projective action of all matrices
in a fixed congruence coset of $M_2(\mathbb Z_p)$. At level $n$ the
answer is a conditional average on one explicitly determined fibre of
the finite projective line. For a primitive character of an unramified
quadratic torus, the formula determines the integral ideal of
differences of the output values in each Smith stratum. Passing from
invertible to primitive singular centres removes exactly one factor
of $\zeta-1$, where $\zeta$ generates the image of the character.
The calculation applies also to cosets that contain singular matrices.
\end{abstract}

\section{The finite projective line and its coset average}

Write $X_n=\mathbb P^1(\mathbb Z/p^n)$ for the primitive pairs modulo
$p^n$ up to multiplication by a unit, where $n\ge1$.
It has $(p+1)p^{n-1}$ elements. Write $X_0$ for a one-point set.
Reduction defines maps $X_n\to X_j$ for $0\le j\le n$.
For a function $f:X_n\to\mathbb C$, denote its average over the
fibre above $x\in X_j$ by $M_j f(x)$. In particular $M_0 f$ is
the average over the entire line and $M_n f=f$.

The projective reduction map $\PP(\Qp)\to X_n$ makes $f$ a locally
constant function on $\PP(\Qp)$. The quotient of these boundary
functions by constants is the smooth Steinberg representation,
with action $(\pi(g)f)(z)=f(g^{-1}z)$. We first compute the operator
on functions, so that the constant term is explicit.

Give $M_2(\Zp)$ additive Haar measure $dX$ of total mass one. For
$a\in M_2(\Zp)$ define
\[
 (\mathcal C_{a,n}f)([z])
 =p^{4n}\int_{a+p^nM_2(\Zp)} f([X^{-1}z])\,dX ,
\]
where $z\in\Zp^2$ is primitive. In an argument of $f$, square
brackets on a nonzero vector denote the projective direction,
reduced modulo $p^n$ after division by the common power of $p$.
The locus $\det X=0$ has additive measure zero. The prefactor makes
the measure on the specified coset a probability measure. In
particular, $\mathcal C_{a,n}$ preserves constants.

\begin{theorem}\label{avg}
Let $v=\adj(a)z$ and put $h=\min(n,v_p(v_1),v_p(v_2))$, using
$v_p(0)=+\infty$. Then
\[
 (\mathcal C_{a,n}f)([z])=
 \begin{cases}
 (M_{n-h} f)([p^{-h}v]\bmod p^{n-h}), & h<n,\\
 M_0 f, &h=n .
 \end{cases}
\]
In particular the output factors through $X_n$.
\end{theorem}

\begin{proof}
In dimension two the adjugate is an invertible integral linear map
of the four matrix entries. For $Y$ Haar-uniform on $M_2(\Zp)$,
$\adj(Y)z$ is Haar-uniform on $\Zp^2$: the map is an additive
surjection, since $z$ has a unit coordinate. For $X=a+p^n Y$ we have
\[
 [X^{-1}z]=[\adj(X)z],
 \qquad \adj(X)z=\adj(a)z+p^n\adj(Y)z
\]
outside the singular locus.

If $h<n$, division by $p^h$ leaves a uniform lift of one fixed
primitive vector modulo $p^{n-h}$. Its projective reductions at
level $n$ are uniform on the claimed fibre: a projective lift has
exactly $p^h$ vector representatives extending that fixed vector.
If $h=n$, division by $p^n$ leaves a Haar-uniform vector on
$\Zp^2$. After projectivization its law is invariant under
$\mathrm{GL}_2(\Zp)$, which is transitive on $X_n$.
It is therefore uniform. The same argument shows that all quantities
on the right depend only on $a\bmod p^n$ and $[z]\in X_n$.
\end{proof}

For $n\ge2$ call $f$ primitive at level $n$ if $M_{n-1}f=0$.
Then $M_jf=0$ for every $j<n$. The formula becomes
\begin{equation}\label{adj}
 (\mathcal C_{a,n}f)([z])=
 \begin{cases}
 f([\adj(a)z]),&\adj(a)z\not\equiv0\pmod p,\\
 0,&\adj(a)z\equiv0\pmod p .
 \end{cases}
\end{equation}
For example, for $a=\operatorname{diag}(1,p^d)$ with $1\le d<n$,
the output on $[t:1]$ equals $f([p^dt:1])$. On the complementary
chart with second coordinate divisible by $p$ it is zero. When
$d\ge n$ this is the indicator of the affine chart multiplied by
$f([0:1])$. This last identity is a coset average even with
$a=\operatorname{diag}(1,0)$.

\section{Integral content on the Smith strata}

Let $p$ be odd, and identify $\Zp^2$ with the ring of integers
$\OO_L$ in the unramified quadratic extension $L/\Qp$. Every
primitive vector is then a unit of $\OO_L$. The group
\[
 T_n=(\OO_L/p^n\OO_L)^\times/(\mathbb Z/p^n)^\times
\]
acts simply transitively on $X_n$ by multiplication. Take $[1]$ as
base point to identify the two sets. Let $\chi:T_n\to\mathbb C^\times$
be a character of $p$-power order nontrivial on
$\ker(T_n\to T_{n-1})$, with $n\ge2$. The function
$f_\chi([t])=\chi(t)$ has zero average on every fibre to $X_{n-1}$
by character orthogonality.

Let $\zeta$ generate the image of $\chi$, fix a place of
$\mathbb Q(\zeta)$ above $p$, and let $\OO$ be its local ring of
integers. For an $\OO$-valued function $g$ define its content in
the Steinberg quotient by
\[
 \cont([g])=(\,g(z)-g(z'):\ z,z'\in X_n\,)\subseteq \OO .
\]
Indeed, after fixing $z_0$, subtraction of $g(z_0)$ identifies
$\OO^{X_n}/\OO\mathbf1$ with the free lattice of functions zero
at $z_0$. This also defines fractional content for functions with
values in the coefficient field.

\begin{corollary}\label{ct}
For the probability coset average we have
\[
 \cont([\mathcal C_{a,n} f_\chi])=
 \begin{cases}
  (\zeta-1),&\operatorname{rank}(a\bmod p)=2,\\
  \OO,&\operatorname{rank}(a\bmod p)=1,\\
  (0),&\operatorname{rank}(a\bmod p)=0.
 \end{cases}
\]
Thus every primitive singular Smith stratum, including determinant
zero modulo $p^n$, attains the unit ideal.
\end{corollary}

\begin{proof}
An invertible centre permutes $X_n$ in \eqref{adj}. Its output
therefore takes all powers of $\zeta$; the ideal of their differences
is $(\zeta-1)$. For a rank-one centre, $\adj(a)\bmod p$ has
one-dimensional kernel and one-dimensional image. On the residue
direction of the kernel the output vanishes, and on every other
residue direction it is a root of unity. The difference of a unit
and zero is a unit. If $a\bmod p=0$ then also $\adj(a)\bmod p=0$,
and the output vanishes everywhere by \eqref{adj}.
\end{proof}

For clarity about scaling, let $\Phi$ run over all integral test
functions on $M_2(\mathbb Z/p^n)$. The additive operator
\[
 z\longmapsto \int_{M_2(\Zp)}\Phi(X)f_\chi([X^{-1}z])\,dX
\]
is a linear combination of the coset averages with common scalar
$p^{-4n}$. Its coefficient differences generate the fractional
ideal $p^{-4n}\OO$ by Corollary~\ref{ct}. Restricting the test
functions to invertible centres gives $p^{-4n}(\zeta-1)\OO$.
Integration here is first taken over the complex-valued character
for each coset indicator. Formula~\eqref{adj} gives coefficient-field
values, and the map extends linearly to integral tests.

In the scaling used in the archived toric contraction, the last
integral is multiplied by $p^{2-n}/(p+1)^2$. Hence its full
content in that convention is
\[
 k_n\OO,\qquad k_n=\frac{p^{\,2-5n}}{(p+1)^2}.
\]
In particular, integral linear combinations of these tests cannot
have a coefficient difference equal to $k_n/(\zeta-1)$.
The proof concerns this explicitly normalized local operator;
application to a global automorphic coefficient requires an
independent identification of its test data and measures.

\section*{Sources}
The local non-split torus and Steinberg test-vector setting
occurs in File--Martin--Pitale \cite{FMP} and Vatsal \cite{Vatsal}.
The finite coset average and the content question arose in
entries 363 and 394--398 of the archived BSD research record
\cite{archive}. The adjugate argument gives a proof in the
boundary model. Literature priority for this finite calculation
remains to be assessed.

\begin{thebibliography}{9}
\bibitem{FMP} D.~File, K.~Martin, and A.~Pitale,
\emph{Test vectors and central $L$-values for $\mathrm{GL}(2)$},
\href{https://arxiv.org/abs/1310.1765}{arXiv:1310.1765}.
\bibitem{Vatsal} V.~Vatsal,
\emph{Test vectors for some ramified representations},
\href{https://arxiv.org/abs/1806.07856}{arXiv:1806.07856}.
\bibitem{archive} \emph{BSD research history},
\href{https://github.com/cosmosgeek0/BSD/blob/64194dff2dd10ea2fcfa9bd78bc123876d98e1f4/bsd-research-history.zip}{numbered source archive}, entries 301, 363, 390, and 394--398.
\end{thebibliography}
\end{document}
